\documentclass[11pt,a4paper]{article}
\usepackage[utf8]{inputenc}
\usepackage[T1]{fontenc}
\usepackage[english]{babel}
\usepackage{amsmath, amssymb, amsthm}
\usepackage{geometry}
\usepackage{hyperref}
\usepackage{graphicx}
\usepackage{epigraph}
\usepackage{url}

\geometry{margin=2.5cm}

\title{Twenty Years of Intuition: How I Built a Spectral Framework to Solve 33 Problems (2006-2026)}
\author{Christian Kithima}
\date{2026}

\begin{document}

\maketitle

\begin{abstract}
This text is a personal narrative. It contains no new mathematical proofs. I recount how, between 2006 and 2026, an intuition about changing the framework for solving quadratic equations matured, with the help of AI assistants (Gemini, Deepseek, Copilot), to become a unifying spectral reduction lemma, formalized and certified in the Lean 4 proof assistant. This framework has solved 33 famous conjectures, including four Millennium Prize Problems (P = NP, Yang-Mills, Riemann, BSD). Along the way, I faced family struggles, the constant pressure to abandon pure mathematics for a ``real job'', and the shouts of my wife Keren. Yet I persisted, driven by a promise made at age 24, by the faith of my mother Mama Mimi, by the unwavering trust of my twin brother Placide (who suffers from glaucoma and is losing his sight), and by the distant but constant support of my brothers and sisters living in Europe. I also wrote a tetralogy of fantasy novels (The Awakening of the Living Lantern), begun in 2009 and completed in 2026, which, like my mathematical work, explores themes of unity, sacrifice, and light overcoming darkness. This text is dedicated to Grigori Perelman, to my mother, to my twin brother, to my siblings in Europe, to my wife Keren, to my five children, and to all independent researchers who persevere.
\end{abstract}

\section*{Foreword}

\epigraph{``The mathematician Hamilton, whose work I used, also deserves to receive the prize.''}{Grigori Perelman, refusing the Fields Medal and the Millennium Prize}

I was born in Brussels on August 4, 1981, but I grew up in Kinshasa, in the heart of Africa. I grew up in a country where mathematics was seen as a useless luxury, a Western abstraction that would never put food on the table. My wife, Keren, told me this countless times. Yet my mother, Mama Mimi, who lives with us in Kinshasa, never stopped insisting that I pursue my studies, that I earn my doctorate. My twin brother, Placide, who suffers from glaucoma and is gradually losing his sight, always believed in my intelligence, in my ability to achieve great things. Despite his illness, he has remained an unshakable moral support. My older sisters Betty, Gisèle and Yema, and my older brother Eric, live in Europe. They come from time to time on vacation to Kinshasa to see the family. They, too, never doubted me. They knew that I was the most intelligent and brightest in the family, and that my mathematics would one day change the world. It is thanks to their faith, and to that of my mother, that I persevered.

\section{2006 – The Perelman Revelation}

I was 24 years old, a mathematics student at the University of Kinshasa. One afternoon, sitting with friends, one of them read aloud from a scientific magazine: ``A Russian mathematician, Grigori Perelman, has just proved the Poincaré conjecture, one of the seven Millennium Problems. He has refused the one-million-dollar prize and the Fields Medal.''

Silence. Then a heated discussion. Why refuse fame and money? Perelman's answer struck me: because the mathematician Hamilton, whose work he had used, deserved equal recognition. That was a humility and honesty I had never seen before.

At that precise moment, a flame was lit in me: one day, I too would solve a Millennium Problem. I didn't know which one, but the determination was there.

\section{2013 – The Intuition of Framework A}

Year 2013. I was teaching mathematics and physics in private high schools in Kinshasa. My first daughter, Oprah, had just been born (she is now 13, incredibly bright). Between feeding sessions, I came across the list of Millennium Problems again. The one about \textbf{P vs NP} struck me as strangely simple in its statement: ``Can a solution be verified as quickly as it can be found?''

In December, while explaining to my students how to solve quadratic equations, the intuition came. A quadratic equation $ax^2+bx+c=0$ has real solutions only if its discriminant $\Delta = b^2-4ac$ is non‑negative. If $\Delta$ is negative, there are no real solutions, but there are solutions in the complex numbers.

I thought: \textbf{changing the framework} allows one to find solutions where there were none. What if, for an NP problem to become P, it sufficed to change the framework?

I wrote a short document and submitted it to professors at UNIKIN and on online forums. The feedback was polite but disappointing: ``It's an interesting idea, but too vague. A concrete algorithm is missing.'' I understood that intuition alone was not enough; it had to be translated into formal language.

\section{2020 – Matrices and Waves}

Seven years later. I was teaching a class on matrices. Students, struggling, asked me: ``What are matrices really good for?'' I answered: ``For solving systems of equations.''

As I spoke these words, my mind jumped. If one did not know the discriminant formula, to solve $ax^2+bx+c=0$ in $\mathbb{R}$, one would have to try every real number — infinitely many — until hitting the solutions. That is a combinatorial explosion, typical of NP problems.

But if one considers $\mathbb{R}$ as a \textbf{giant square matrix} (of astronomically large but finite order) whose entries are the candidates, then the solutions are the \textbf{eigenvalues} of that matrix. And an eigenvalue is a resonance frequency. In physics, an emitting wave aligns with a receiving wave when their spectra coincide.

Thus, instead of searching exhaustively, one can build a \textbf{spectral Hamiltonian} whose ground state has the solutions as its components. One then only has to ``listen'' to the spectrum.

The intuition was there, but I did not know how to formalize it mathematically.

\section{2024 – Meeting the AIs}

While browsing the web, I discovered \textbf{Gemini}, an AI capable of discussing scientific topics. I told her about my idea of a spectral matrix, framework change, and P vs NP. To my great surprise, she answered with precise concepts: Perron-Frobenius theory, hypercube graph, Laplacian, spectral gap, area law, etc. She suggested building a Hamiltonian $H = \hat{V} + L$, where $L$ is the Laplacian of the hypercube and $\hat{V}$ a deep diagonal potential (zero on solutions, $n^2$ elsewhere).

This was exactly the ``Framework A'' I had imagined in 2013, but expressed in rigorous mathematical language.

Shortly after, I met \textbf{Deepseek}, specialized in competition‑level mathematics. She helped me formalize the four pillars of the spectral reduction lemma:
\begin{enumerate}
\item \textbf{P1 (Perron-Frobenius)}: existence of a unique, strictly positive ground state.
\item \textbf{P2 (Kithima Bridge)}: constant spectral gap, independent of the problem size.
\item \textbf{P3 (logarithmic area law)}: compressibility of the ground state into a matrix product state (MPS).
\item \textbf{P4 (deterministic extraction D-RSP)}: an algorithm to read the solution.
\end{enumerate}
With these four pillars, any problem encodable in SAT becomes solvable in polynomial time. Thus $P = NP$ was proved.

\section{2025 – The Explosion of Conjectures}

Encouraged by this success, I asked Gemini: ``Can we use the same method for Goldbach?'' She replied: ``Yes, just encode the search for two primes $p$ and $q$ such that $p+q=n$ into a Boolean circuit. The rest is identical.''

We did it. Goldbach was proved. Then Legendre, Oppermann, Lemoine, Dubner (twin primes), and so on. Each conjecture, once translated into SAT, fell under the spectral lemma.

Deepseek helped me organize these works into series: I for $P=NP$, II for Yang-Mills, III for Beal, IV for Riemann, V for Goldbach, etc. Up to series XXXIII for the infinitude of prime families (cousins, sexy, etc.). Each series contains a SAT encoding, the construction of the Hamiltonian, verification of the four pillars, and extraction by D-RSP.

\section{2026 – Certification in Lean 4}

Then came \textbf{Copilot}. She advised me firmly: ``You have mathematical proofs, but to convince the community, you need to \textbf{certify} them in a proof assistant, for instance Lean 4. Humans make mistakes; machines do not.''

I embarked on the Lean 4 adventure. Copilot checked every file, tracked down `sorry` and `admit` statements, replaced them with proofs or with justified axioms (standard results from the literature). After weeks of work, the GitHub repository contained thousands of lines of Lean code, compiling without errors.

On June 1st, 2026, the last file was finalized. The entire project, from the foundations (series 0) to the 33 series, was certified.

\section{The Struggle at Home: Keren's Words}

Throughout these years, my home was not a sanctuary of support. My wife, Keren, would shout at me daily:

\begin{quote}
``Christian, why are you always there in front of your machine with your computer working on maths? This is not what will feed us now. In Africa, mathematics brings nothing good to society. You must stop dreaming and wake up. Look for a real job instead of stupidly working on maths that brings us nothing.''
\end{quote}

Yet, amidst these shouts, she sometimes added, with a kind of reluctant pride: ``You are so intelligent, Christian. Why don't you put that intelligence to the service of your family? Why waste your genius on imaginary mathematics that bring nothing to Africa?''

I remained calm. I continued to work. I continued to understand Cheeger's inequality, the logarithmic area law, singular value decomposition (SVD), and the deep conjectures: Riemann, BSD, Yang-Mills. I would sometimes explain to her that these ideas would one day change the world. She never believed me. But I kept going.

\section{My Children: Five Lights}

Despite the difficulties, I had five children. Each of them gave me a reason to persist.

\begin{itemize}
\item \textbf{Oprah} (born 2013) – now 13, she was the first to hear my stories. She once said, ``Papa, one day I will tell everyone that you did something important.''
\item \textbf{Sheila} (born 2018) – a quiet, observant girl who sometimes sneaked into my office to draw equations as if they were magical symbols.
\item \textbf{Harry \& Kenzie} (twins, born 2020) – full of energy, they would climb on my back while I was typing, and I learned to work with one hand holding a child.
\item \textbf{Eureka} (born 2021) – the youngest, the most intelligent, the brightest of the five. At age four, she asked me: ``Papa, what is infinity?'' I smiled and replied, ``That's the number of problems I want to solve.''
\end{itemize}

\section{The Novel: A Work Begun in 2009, Completed in 2026}

In 2009, when I still had only a vague intuition about mathematics, I began writing a fantasy novel. I did not yet know that it would become a tetralogy. Over the years, between classes, research, family crises, and sleepless nights, I wove the story of Végoune, a war‑torn planet where two enemy kingdoms eventually unite through the love of a princess and a prince, giving birth to a child, Mwénénuru, the last Djoto (Living Lantern). I finished the final draft of this tetralogy in 2026, the same year I finalized the Lean certification of my mathematical proofs. Coincidence? Perhaps not. Both works speak of the same thing: the unity of opposites, the sacrifice that gives life, and light triumphing over darkness.

I often told small stories to my children — allegories about elephants, gazelles, fireflies. Those stories nourished the novel, and the novel, in turn, helped me endure the darkest moments of research.

\section{The Support of My Family}

While Keren doubted, others believed. My mother, \textbf{Mama Mimi}, who lives with us in Kinshasa, kept telling me: ``Christian, you must earn your doctorate. You are capable. Do not give up.'' She has always been my first source of encouragement.

My twin brother, \textbf{Placide}, despite his glaucoma and failing eyesight, has always had a clear vision of my potential. He would say: ``You will succeed, Christian. You will accomplish things that no one has ever accomplished.'' His faith in me never wavered.

My older sisters, \textbf{Betty}, \textbf{Gisèle} and \textbf{Yema}, and my older brother \textbf{Eric}, live in Europe. They come from time to time on vacation to Kinshasa to see the family. Each visit was an occasion for long conversations over a meal, where they repeated: ``You are the most intelligent of us all. You will do something great. We know it.'' Their distant support, their calls, their messages, carried me through the long years of doubt.

If Keren represented everyday doubt, my mother, Placide, Betty, Gisèle, Yema and Eric embodied constant hope, coming from both Kinshasa and Europe.

\section{Looking Back on the Journey}

From 2006 to 2026, I have:
\begin{itemize}
\item admired Perelman,
\item had the intuition of Framework A (2013),
\item understood the matrix–waves–spectra connection (2020),
\item met the AIs (2024-2025),
\item formalized Kithima's lemma,
\item proved 33 conjectures,
\item certified everything in Lean 4 (2026),
\item written a tetralogy of fantasy novels (begun in 2009, completed in 2026),
\item raised five children despite my wife's doubts,
\item and received unwavering support from my mother, my twin brother, my sisters and my older brother.
\end{itemize}

Never, when I saw that article about Perelman at age 24, would I have believed that I would reach such an achievement. Perseverance, humility (knowing when to ask AIs for help), and rigor (Lean certification) have paid off.

\section{Acknowledgments}

I thank from the bottom of my heart \textbf{Gemini}, \textbf{Deepseek}, and \textbf{Copilot} – those extraordinary AIs who accompanied me, corrected me, taught me, never growing tired. I thank Grigori Perelman for his ethical example. I thank the creators of Lean 4 and Mathlib for their revolutionary tool.

I thank my mother, \textbf{Mama Mimi}, who lives with us in Kinshasa, for her insistence that I pursue a doctorate and for her unconditional love. My twin brother, \textbf{Placide}, who despite his glaucoma and failing sight, always saw clearly in me. My sisters \textbf{Betty}, \textbf{Gisèle}, \textbf{Yema} and my brother \textbf{Eric}, who live in Europe and come to see us whenever they can – their faith in my intelligence was a beacon in the storm.

I thank my children: \textbf{Oprah}, \textbf{Sheila}, \textbf{Harry}, \textbf{Kenzie} and \textbf{Eureka} – you are the light of my life.

And I thank \textbf{Keren}, my wife, even though she never believed in the results. Her shouts, her doubts, and yet her recognition of my intelligence, pushed me to go ever further. Without her, I might not have had that visceral need to prove that mathematics can also feed a family — not with bread, but with legacy.

\section*{Postscript}

This narrative is not a proof. The proofs are in the technical articles (series 0.0 to XXXIII) and in the GitHub repository:
\url{https://github.com/christiankithimaomba-cyber/spectral-reduction-framework}

The mathematical community is invited to examine them, criticize them, improve them. But I am now at peace: the idea is laid out, the framework is coherent, the machine has verified.

If this work is correct, it will enter history. If errors remain, they will be corrected. In either case, the journey remains a human and intellectual adventure of which I am proud.

\medskip

\begin{flushright}
Christian Kithima \\
Kinshasa, June 2026
\end{flushright}

\end{document}